\documentclass[10pt,a4paper]{amsart}
\usepackage[margin=19mm]{geometry}
\usepackage{amsmath,amssymb,mathtools}
\usepackage[scr=rsfs,cal=euler]{mathalfa}
\usepackage{xfrac}
\usepackage[unicode,hidelinks,bookmarks=false]{hyperref}
\newcommand{\ie}{i.e.\ }
\newcommand{\cf}{cf.\ }
\newcommand{\resp}{respectively\ }
\newcommand{\Subsection}{\subsection}
\providecommand{\nobreakdash}{\nobreak\hbox{-}}
\setlength{\emergencystretch}{2em}
\multlinegap0pt
\usepackage{bm}
\usepackage{bbm}
\usepackage{dsfont}
\usepackage{xspace}
\usepackage{cancel}
\usepackage{mathtools}
\usepackage{amsthm}
\makeatletter
\@ifundefined{theorem}{\newtheorem{theorem}{Theorem}}{}
\@ifundefined{observation}{\newtheorem{observation}[theorem]{Observation}}{}
\@ifundefined{proposition}{\newtheorem{proposition}[theorem]{Proposition}}{}
\makeatother
\newcommand\abs[1]{\left|#1\right|}
\newcommand\dist[2]{\operatorname{dist}\left(#1,#2\right)}
\renewcommand\deg[1]{\operatorname{deg}\left(#1\right)}
\title[On the forts and related parameters of the hypercube graph]{On the forts and related parameters of the hypercube graph}
\author{Boris Brimkov, Thomas R. Cameron, Owen Grubbs}
\date{Source: arXiv:2507.10826v1 (CC BY 4.0)}
\begin{document}
\maketitle

\noindent\emph{Source and licence.} On the forts and related parameters of the hypercube graph. Version arXiv:2507.10826v1 (CC BY 4.0), distributed under the Creative Commons Attribution 4.0 International licence, as recorded in the arXiv licence metadata for this exact version and shown on the abstract page https://arxiv.org/abs/2507.10826v1. Full rights notice: licenses/arxiv-2507.10826-CC-BY-4.0.txt.\par\smallskip

\noindent\textit{Editorial scope.} Counted unit: the Proposition whose source label is \texttt{prop:max\_dist4} in the article (source0.tex lines 367--371, source label \texttt{prop:max\_dist4}), delivered together with the article's own complete original proof of it (source0.tex lines 372--404, one \texttt{proof} environment of 33 lines). No mathematics is written, repaired or re-derived: statement and proof are the article's own text line for line. Required context retained uncounted: obs:dreg (lines 164--166); obs:neighbors (lines 184--187); prop:4cycle (lines 208--211); prop:max\_dist6 (lines 333--336). Every definition, display and label of those lines is retained unchanged. Omitted from this package and not redistributed: 1--163, 167--183, 188--207, 212--332, 337--366, 405--1254. Nothing inside a retained excerpt is omitted or altered: every retained body is delivered exactly as the article's lines are written, so no omission receipt of any kind accompanies this package. Citations: the retained excerpts carry no citation command at all, so no bibliography entry of the article is transcribed and no reference is redistributed. Reference adaptation: none was needed and none was made; every reference inside the delivered unit resolves inside it, so no unresolved reference or citation is produced. Printed slips kept literal: no printed word, symbol, hypothesis or number of the counted statement or of its proof was altered. Layout and class: the article's own class is \texttt{article} with packages including amssymb, amsmath, mathrsfs, amsthm, enumerate, graphicx, color, geometry, float, caption, subcaption, xcolor, tikz, optidef; the wrapper is the lane's pinned \texttt{amsart} editorial wrapper at 10pt on A4, which restates the theorem-like environments the retained text needs and adds the article's own macro definitions that the retained text needs (\texttt{\textbackslash abs}, \texttt{\textbackslash deg}, \texttt{\textbackslash dist}), with extra wrapper packages (bm, bbm, dsfont, xspace, cancel, mathtools). The delivered numbering is the unit's own. Assets: no retained span contains a figure environment, a graphics-inclusion command, a PDF-page inclusion, a table or a TikZ picture; no image file is copied into this package, the package declares no asset dependency and no image file is redistributed with it. Licence: version arXiv:2507.10826v1 of this article is distributed under the Creative Commons Attribution 4.0 International licence, as recorded in the arXiv licence metadata for this exact version and shown on the abstract page https://arxiv.org/abs/2507.10826v1. A full rights notice is delivered at licenses/arxiv-2507.10826-CC-BY-4.0.txt. Characters: every retained excerpt is pure ASCII, so the delivered engine, XeLaTeX with its default Latin Modern text font, reports no missing character.
\begin{observation}\label{obs:dreg}
For all $v\in V(Q_{d})$, $\deg{v}=d$. 
\end{observation}

\begin{observation}\label{obs:neighbors}
Let $u\in V(Q_{d})$.
For all $v,w\in N(u)$, $v\notin N(w)$. 
\end{observation}

\begin{proposition}\label{prop:4cycle}
Let $u,v\in V(Q_{d})$.
If $N(u)\cap N(v)\neq\emptyset$, then $\abs{N(u)\cap N(v)}=2$.
\end{proposition}

\begin{proposition}\label{prop:max_dist6}
Let $d\geq 6$ and let $u,v\in V(Q_{d})$ such that $\dist{u}{v}\geq 6$.
Then, there is no minimum fort of $Q_{d}$ that contains both $u$ and $v$. 
\end{proposition}

\begin{proposition}\label{prop:max_dist4}
Let $d\geq 4$ and let $u,v\in V(Q_{d})$ such that $\dist{u}{v}=4$.
If $d=4$, then any fort $F$ containing $u$ and $v$ satisfies $\abs{F}\geq d$.
If $d\geq 5$, then any fort $F$ containing $u$ and $v$ satisfies $\abs{F}>d$. 
\end{proposition}

\begin{proof}
Let $F=\{u,v\}\subset V(Q_{d})$.
Note that $F$ is not a fort of $Q_{d}$ since every $w\in N(u)\cup N(v)$ satisfies $\abs{N(w)\cap F}=1$. 
Hence, by Observation~\ref{obs:dreg}, there are $\abs{N(u)\cup N(v)}=2d$ fort violations of $F$.
There are two options to repair these fort violations by adding a vertex to $F$: Either we add a vertex from $N(u)\cup N(v)$, or we add a vertex that is not in $N[u]\cup N[v]$ but is adjacent to a vertex in $N[u]\cup N[v]$.
Obviously, by adding a vertex to $F$ we may introduce new fort violations, but we will ignore this effect for now.

Let $u'\in N(u)$.
By Observation~\ref{obs:neighbors}, $u'$ is not adjacent to any vertex in $N(u)$.
Furthermore, since $\dist{u}{v}=4$, $u'$ is not adjacent to any vertex in $N(v)$. 
Therefore, adding $u'$ to $F$ resolves exactly $1$ fort violation of $F$, namely $u'$.
Similarly, adding $v'\in N(v)$ to $F$ resolves exactly $1$ fort violation of $F$, namely $v'$. 

Let $p\in V(Q_{d})\setminus\left(N[u]\cup N[v]\right)$.
Then, by Proposition~\ref{prop:4cycle}, $p$ may be adjacent to no vertices in $N(u)\cup N(v)$, exactly $2$ vertices in $N(u)$ or $2$ vertices in $N(v)$, or exactly $2$ vertices in $N(u)$ and $2$ vertices in $N(v)$. 
In the latter case, $p$ is adjacent to $u',u''\in N(u)$ and $v',v''\in N(v)$ such that $d(u',v)=d(u'',v)=3$ and $d(v',u)=d(v'',u)=3$.
There are exactly $4$ vertices in $N(u)$ a distance of $3$ away from $v$; similarly, there are exactly $4$ vertices in $N(v)$ a distance of $3$ away from $u$.
Moreover, all $8$ of these fort violations can be resolved by adding $p_{1},p_{2}\in V(Q_{d})\setminus\left(N[u]\cup N[v]\right)$ to $F$, where $p_{1}$ and $p_{2}$ are a distance of $2$ away from both $u$ and $v$ and $\dist{p_{1}}{p_{2}}=4$.
After adding $p_{1}$ and $p_{2}$ to $F$, we still have $2d-8$ fort violations from $N(u)\cup N(v)$.
Furthermore, every vertex that can be added at this point resolves at most $2$ of these fort violations.
Hence, we must add at least $d-4$ more vertices to $F$ in order to resolve all fort violations.

Suppose that $d=5$.
Then, there are exactly two fort violations that remain from $N(u)\cup N(v)$, which we denote by $u'\in N(u)$ and $v'\in N(v)$.
Since $d(u',v)=d(v',u)=5$, these fort violations cannot be resolved by a single vertex. 
Hence, when $d=5$, at least $4$ vertices must be added to $F$ in order resolve all fort violations. 

Finally, suppose that $d\geq 6$. 
Then, there are at least two fort violations that remain from $N(u)$, which we denote by $u',u''$.
Note that $d(u',v)=d(u'',v)=5$. 
Moreover, by Proposition~\ref{prop:4cycle}, there exists a unique $w$ such that $N(w)\cap N(u)=\{u',u''\}$; hence, we can resolve two fort violations with $w$.
However, $d(w,v)=6$ and Proposition~\ref{prop:max_dist6} implies that no minimum fort of $Q_{d}$ will contain both $w$ and $v$.
\end{proof}

\end{document}
